\chapter{Адаптивная кост модель}\label{ch:ch1}

\section{Краткое описание предлагаемой модели}\label{sec:ch1/sec1}

Предлагаемая модель по-разному устанавливает параметры стоимостной модели, связанные с операциями ввода-вывода диска (\textit{I/O cost}),  и параметры, относящиеся к загрузке процессора (\textit{CPU cost}). Для их корректной оценки требуются разные подмодели, поскольку CPU-параметры не зависят от текущего состояния базы данных, тогда как параметры, связанные с диском, могут изменяться в зависимости от состояния буферного кэша базы данных, которое, в свою очередь, меняется практически после каждого выполненного запроса.

Для обеспечения быстрого выполнения запросов модель вычисляет оптимальные I/O-параметры непосредственно перед запуском каждого запроса. Для этого оценивается объём запрашиваемых данных, уже находящихся в буферном кэше СУБД, и прогнозируется объём данных, который предстоит прочитать с диска. Для построения точных прогнозов этих величин модель использует исторические данные о выполнении предыдущих запросов.


Для вычисления оптимальных CPU-параметров модель непрерывно отслеживает статистику исполнения операторов плана и подбирает параметры так, чтобы оценка времени выполнения (полученная на этапе планирования) как можно ближе соответствовала фактическому времени выполнения (полученному после завершения запроса). Таким образом, ACM непрерывно корректирует параметры стоимостной модели, повышая точность оценок стоимости  планов и позволяя оптимизатору (планировщику) выбирать оптимальный план выполнения запроса.

В следующем разделе подробно описываются двух компонентов Адаптивной стоимостной модели: сначала модели, связанной с дисковой подсистемой, затем - модели, связанной с CPU. Стоит отметить, что одна единица стоимости представляет собой условную величину, обозначающую время, необходимое для обработки одной страницы при последовательном сканировании. Все остальные значения выражаются как доли этой базовой величины. Например, \(cpu\_tuple\_cost\) по умолчанию равен 0.01, что означает , что обработка 100 строк требует столько же времени, сколько последовательное чтение одной страницы с диска.  Значение \(cpu\_index\_tuple\_cost\) установлено на 0.005, а \(cpu\_operator\_cost\) - на 0.0025. В свою очередь, \(random\_page\_cost\) равен 4, что указывает: время, необходимое для случайного доступа к странице, превышает время последовательного доступа в 4 раза.

Следовательно, чтобы увеличить стоимость последовательного сканирования страницы, можно уменьшить значения остальных четырех параметров; и наоборот, чтобы увеличить стоимость всех операций, кроме \(seq\_page\_cost\), достаточно уменьшит только \(seq\_page\_cost\). Таким образом, для настройки дисковых и зависящих от CPU параметров требуется использовать лишь эти четыре параметра. В данном исследовании мы фиксируем \(seq\_page\_cost\), так как в стандартной конфигурации он равен 1, что позволяет сравнивать изменения прочих параметров относительно их предыдущих значений.


\section{Дисковые параметры}\label{sec:ch1/sec2}

Стандартная стоимостная модель различает стоимость последовательного доступа к странице и стоимость случайного доступа к странице за счёт различий в параметрах стоимости. Однако такой подход не позволяет оценить, какая часть запрошенных страниц уже находится в буфере и, следовательно, не требует ни последовательного, ни случайного чтения с диска. На этапе планирования СУБД не знает, какая доля запрошенных данных уже находится в памяти, что неизбежно приводит к ошибкам в оценке стоимости выполнения. Предлагаемая система повышает точность оценивания стоимости за счёт динамической корректировки стоимости случайного дискового доступа на основе информации о доступности данных в буфере.

При выполнении запроса база данных собирает статистику использования данных из буфера (buffer hit) и с диска (read). Эта связь важна, так как показывает, какая часть запрошенных из таблицы данных находится в памяти. Прежде всего, предлагаемая модель вычисляет параметр \(random\_page\_cost\) отдельно для каждой таблицы, поскольку одна таблица может полностью находиться в буферном кэше, а другая отсутствовать. Если вся таблица находится в буферном кэше, разумно приравнять \(random\_page\_cost\) к \(seq\_page\_cost\), указывая оптимизатору, что время чтения будет эквивалентным. Таким образом, при расчете стоимости система использует значение \(random\_page\_cost\), соответствующее сканируемой таблице.

Рассмотрим, как динамически вычисляется параметр \(random\_page\_cost\) для таблицы. Как отмечено выше, точная оценка \(random\_page\_cost\) зависит от того, какая доля запрошенных данных будет взята из буфера, а какая - прочитана с диска. Это можно выразить показателем «hit ratio»:

\begin{equation}
    R(q) = \frac{hit}{hit + read}
\end{equation}
где
\begin{itemize}
    \item \(q\) - запрос, для которого вычисляется показатель;
    \item \(hit\) - количество страниц таблицы, использованных в запросе и прочитанных из памяти;
    \item \(read\) - количество страниц таблицы, использованных в запросе и прочитанных с диска.
\end{itemize}

Однако эта информация становится известна лишь после завершения запроса. Тем не менее можно попытаться оценить данный показатель, используя статистику выполнения предыдущих запросов. Для каждой таблицы система накапливает статистику «hit ratio»:

\begin{equation}
    R(q_{t}), R(q_{t-1}), R(q_{t-2}), ...
\end{equation}
где
\begin{itemize}
    \item \(R(q_{t})\) - hit ratio запроса \(q_{t}\) выполненного в момент времени \(t\).
\end{itemize}

Соответственно, один из способов оценить hit ratio \(R(q_{t+1})\) для нового запроса \(q_{t+1}\) до его выполнения — построить модель, которая предсказывает hit ratio по предыдущим значениям \(R(q_{t})\), \(R(q_{t-1})\), ... . Здесь можно использовать различные модели прогнозирования временных рядов (например, ARIMA, экспоненциальное сглаживание), однако мы предлагаем более простую модель, основанную лишь на предыдущем hit ratio и коэффициенте деградации, поскольку эта простая модель показала результаты не хуже более сложных:
\begin{equation}
    R(q_{t+1}) = R(q_{t}) \cdot D ,
\end{equation}
\begin{equation}
    D = \frac{1 + (qc - tc)}{1 + (qc - tc)^2} 
\end{equation}
где
\begin{itemize}
    \item \(D\) - коэффициент деградации;
    \item \(qc\) - счётчик запросов, текущее общее число обращений ко всем таблицам; увеличивается на единицу при каждом чтении любой таблицы;
    \item \(tc\) - счётчик, фиксирующий момент последнего доступа к целевой таблице; при обращении к целевой таблице текущее значение \(qc\) записывается в \(tc\).    
\end{itemize}

Поясним смысл коэффициента деградации. Важно учитывать, что после запроса к таблице \(T_{1}\) могут выполняться другие запросы к таблицам \(T_{2}, ..., T_{n}\), которые способны вытеснить данные таблицы \(T_{1}\) из буфера. Чем больше запросов к другим таблицам было выполнено после последнего доступа к \(T_{1}\) тем меньше данных \(T_{1}\) останется в буфере. Коэффициент деградации позволяет учесть этот эффект.

Далее, необходимо связать \(R(q_{t+1})\) с \(random\_page\_cost\) таким образом, чтобы при максимальном прогнозном hit ratio значение \(random\_page\_cost\) совпадало с \(seq\_page\_cost\). Следовательно:
\begin{equation}
    \begin{split}
        random\_page\_cost_{new} = random\_page\_cost_{def} \cdot \\ (1 - R(q_{t+1})) +  seq\_page\_cost \cdot R(q_{t+1})
    \end{split}
\end{equation}
где
\begin{itemize}
    \item \(random\_page\_cost_{def}\) - значение по умолчанию для \(random\_page\_cost\),
    \item \(random\_page\_cost_{new}\) - новое значение для \(random\_page\_cost\), которое будет использовано при расчете стоимости сканирования конкретной таблицы.
\end{itemize}

\begin{figure}[ht]
    \centerfloat{
        \includegraphics[scale=0.45]{disk model.PNG}
    }
    \caption{Архитектура дисковой модели.}
    \label{fig:disk_architecture}
\end{figure}

На рисунке ~\cref{fig:disk_architecture} показана схема получения статистики использования буфера при выполнении запросов. Чем больше страниц таблицы читается из памяти, тем ниже стоимость последующих случайных обращений, поскольку большинство страниц может быть обслужено из буфера. Следовательно, стоимость случайного доступа к странице необходимо корректировать для каждой таблицы данных.

Для каждой таблицы количество страниц, прочитанных с диска и из памяти, накапливается до тех пор, пока объём данных не достигнет уровня, достаточного для прогнозирования hit ratio следующего обращения к этой таблице. Затем на этапе планирования стоимость случайного доступа к странице предсказывается с использованием коэффициента деградации, и полученные стоимости «инъектируются» в оптимизатор. Система также учитывает временной интервал между последним обращением к таблице и планируемым, корректируя число попаданий в формуле коэффициента деградации. После внедрения стоимостей планирование запроса продолжается в обычном режиме. По завершении выполнения запроса с «инъецированными» стоимостями новые данные записываются для всех таблиц, к которым был доступ.

\section{CPU параметры}\label{sec:ch1/sec3}


\begin{figure}[h]
    \centering
    \includegraphics[scale=0.5]{cpu model.PNG}
    \caption{Архитектура CPU модели}
    \label{fig:cpu-architecture}
\end{figure}

Другие три параметра стандартной стоимостной модели отвечают за операции,
выполняемые процессором: стоимость обработки одного кортежа (tuples cost),
выполнения одной операции (operator cost) и обработки одного кортежа индекса
(index tuple cost). Стандартные значения этих параметров не отражают специфику
различных аппаратных платформ и конфигураций баз данных, что приводит к неточным
оценкам стоимости запросов. Кроме того, стандартная модель не учитывает
особенности стоимости для разных операторов с индивидуальными формулами расчета
стоимости. Предлагаемая система предоставляет метод сбора статистики и
динамической настройки параметров с учетом конфигурации базы данных и аппаратной
платформы отдельно для каждого оператора плана. На рисунке ~\cref{fig:cpu-architecture} показана схема получения статистики использования CPU во время
выполнения запросов.

После выполнения запроса система собирает ключевую статистику исполнения, такую
как время выполнения операторов, число обработанных кортежей в каждом операторе
и другие показатели. Индивидуальная статистика для каждого типа оператора
накапливается отдельно. На основе этих данных минимизируется расхождение между
значениями, выдаваемыми формулами стоимости операторов, и фактическими временами
их выполнения. Целевая функция строится из легковесных линейных моделей, что не
снижает скорость работы базы данных.

Функция параметров стоимости оператора:

\begin{equation}
    \begin{split}
        & F(c_t, c_o, c_i) = c_t \cdot n_t+c_o \cdot n_o + c_i \cdot n_i + s, \\
        & s = c_s \cdot n_s + c_r \cdot n_r
    \end{split}
    \label{eq:to_min}
\end{equation}
где
\begin{itemize}
    \item \(c_t\) - стоимость обработки одного кортежа,
    \item \(n_t\) - число обработанных кортежей,
    \item \(c_o\) - стоимость выполнения одной операции,
    \item \(n_o\) - число выполненных операций,
    \item \(c_i\) - стоимость обработки одного кортежа индекса,
    \item \(n_i\) - число кортежей индекса,
    \item \(s\) -  стоимость последовательных и случайных чтений для узла при параметрах GUC, полученных из дисковой модели.
\end{itemize}

Исходя из формулы функции стоимости, можно построить простую линейную модель:
\begin{equation}
    L(c,x) = x^T c
    \label{eq:model}
\end{equation}
где вектор \(c=(c_t, c_o, c_i, 1)\) представляет параметры модели, а \(x=(n_t, n_o, n_i, s)\) — независимые переменные. Заметим, что \(s\) рассматривается как независимая переменная, поскольку его значение задаётся дисковой моделью.
Рассмотрим выборку наблюдений \(X\) значений переменных, собранных системой в
процессе выполнения запросов:
\begin{equation}
    X=
    \begin{pmatrix}
        n_t^1 & n_o^1 & n_i^1 & s^1\\
        n_t^2 & n_o^2 & n_i^2 & s^2\\
        ... & ... & ... & ... \\
        n_t^n & n_o^n & n_i^n & s^n\\
    \end{pmatrix}
\end{equation}
Тогда модель \ref{eq:model} можно переписать в матричном виде:
\begin{equation}
    L(c,X) = Xc
    \label{eq:model_matrix}
\end{equation}

Таким образом, задача нахождения оптимальных параметров стоимости сводится к
задаче линейной регрессии:
\begin{equation}
    L(c,X) \rightarrow time \cdot scale\_factor
\end{equation}
где
\begin{itemize}
    \item
    \(time=
    \begin{pmatrix}
        t^1\\
        t^2\\
        ...\\
        t^n\\
    \end{pmatrix}
    \) — собранные времена выполнения оператора,
    \item \(scale\_factor\) — константа, представляющая коэффициент масштабирования
    между временем выполнения и стоимостью.
\end{itemize}
Для оценки параметров модели можно использовать различные методы. В данной работе
использован классический метод наименьших квадратов.

Таким образом, описанная выше процедура позволяет получать точные оценки
параметров стоимости для каждого оператора и устранять ошибки, связанные с
использованием одинаковых параметров стоимости для всех типов операторов.

После сбора статистики и вычисления параметров стоимости система сохраняет
результаты для последующего использования при прогнозировании параметров стоимости
для будущих запросов. При динамической подстройке параметров важно своевременно
реагировать на изменения конфигурации базы данных и другие факторы, влияющие на
стоимости. Поэтому для оперативной реакции на такие изменения система применяет
методы анализа временных рядов для оценки будущих параметров стоимости, когда они
настраиваются индивидуально для каждого типа оператора.
Например, для параметра \(cpu\_tuple\_cost\), предсказанного для моментов времени
\(t_0,..t_{n-1}\): [\(cpu\_tuple\_cost_0\), ..., \(cpu\_tuple\_cost_{n-1}\)],
предсказывается \(cpu\_tuple\_cost_n\) с использованием моделей прогнозирования
временных рядов (например, ARIMA и т. п.). В данной работе использовано
экспоненциальное скользящее среднее:
\begin{equation}
    c_{pred, n} = (1 - \alpha) \cdot c_{n-1} + \alpha \cdot c_{pred, n - 1}
\end{equation}
где
\begin{itemize}
    \item \(c_{pred, i}\) — предсказанное значение параметра стоимости на шаге \(i\),
    \item \(c_i\) — значение параметра стоимости, полученное из линейной модели
    (\ref{eq:model_matrix}) на шаге \(i\),
    \item \(\alpha\) — коэффициент сглаживания.
\end{itemize}

В целом система использует статистику выполнения запросов для калибровки пяти
параметров стоимости, что позволяет планировщику точнее оценивать стоимость для
каждого задействованного типа оператора. Информация для системы включает,
помимо прочего, данные о состоянии буфера, время выполнения запроса и его узлов,
число произведённых кортежей и другие ключевые характеристики.



Мы можем сделать \textbf{жирный текст} и \textit{курсив}.

Сошлёмся на библиографию.
Одна ссылка: \cite[с.~54]{Sokolov}\cite[с.~36]{Gaidaenko}.
Две ссылки: \cite{Sokolov,Gaidaenko}.
Ссылка на собственные работы: \cite{vakbib1, confbib2}.
Много ссылок: %\cite[с.~54]{Lermontov,Management,Borozda} % такой «фокус»
%вызывает biblatex warning относительно опции sortcites, потому что неясно, к
%какому источнику относится уточнение о страницах, а bibtex об этой проблеме
%даже не предупреждает
\cite{Lermontov, Management, Borozda, Marketing, Constitution, FamilyCode,
    Gost.7.0.53, Razumovski, Lagkueva, Pokrovski, Methodology, Berestova,
    Kriger}%
\ifnumequal{\value{bibliosel}}{0}{% Примеры для bibtex8
    \cite{Sirotko, Lukina, Encyclopedia, Nasirova}%
}{% Примеры для biblatex через движок biber
    \cite{Sirotko2, Lukina2, Encyclopedia2, Nasirova2}%
}%
.
И~ещё немного ссылок:~\cite{Article,Book,Booklet,Conference,Inbook,Incollection,Manual,Mastersthesis,
    Misc,Phdthesis,Proceedings,Techreport,Unpublished}
% Следует обратить внимание, что пробел после запятой внутри \cite{}
% обрабатывается ожидаемо, а пробел перед запятой, может вызывать проблемы при
% обработке ссылок.
\cite{medvedev2006jelektronnye, CEAT:CEAT581, doi:10.1080/01932691.2010.513279,
    Gosele1999161,Li2007StressAnalysis, Shoji199895, test:eisner-sample,
    test:eisner-sample-shorted, AB_patent_Pomerantz_1968, iofis_patent1960}%
\ifnumequal{\value{bibliosel}}{0}{% Примеры для bibtex8
}{% Примеры для biblatex через движок biber
    \cite{patent2h, patent3h, patent2}%
}%
.

\ifnumequal{\value{bibliosel}}{0}{% Примеры для bibtex8
Попытка реализовать несколько ссылок на конкретные страницы
для \texttt{bibtex} реализации библиографии:
[\citenum{Sokolov}, с.~54; \citenum{Gaidaenko}, с.~36].
}{% Примеры для biblatex через движок biber
Несколько источников (мультицитата):
% Тут специально написано по-разному тире, для демонстрации, что
% применение специальных тире в настоящий момент в biblatex приводит к непоказу
% "с.".
\cites[vii--x, 5, 7]{Sokolov}[v"--~x, 25, 526]{Gaidaenko}[vii--x, 5, 7]{Techreport},
работает только в \texttt{biblatex} реализации библиографии.
}%

Ссылки на собственные работы:~\cite{vakbib1, confbib1}.

Сошлёмся на приложения: Приложение~\cref{app:A}, Приложение~\cref{app:B2}.

Сошлёмся на формулу: формула~\cref{eq:equation1}.

Сошлёмся на изображение: рисунок~\cref{fig:knuth}.

Стандартной практикой является добавление к ссылкам префикса, характеризующего тип элемента.
Это не является строгим требованием, но~позволяет лучше ориентироваться в документах большого размера.
Например, для ссылок на~рисунки используется префикс \textit{fig},
для ссылки на~таблицу "--- \textit{tab}.

В таблице \cref{tab:tab_pref} приложения~\cref{app:B4} приведён список рекомендуемых
к использованию стандартных префиксов.

В некоторых ситуациях возникает необходимость отойти от требований ГОСТ по оформлению ссылок на
литературу.
В таком случае можно воспользоваться дополнительными опциями пакета \verb+biblatex+.

Например, в ссылке на книгу~\cite{sobenin_kdv} использование опции \verb+maxnames=4+ позволяет
вывести имена всех четырёх авторов.
По ГОСТ имена последних трёх авторов опускаются.

Кроме того, часто возникают проблемы с транслитерованными инициалами. Некоторые буквы русского
алфавита по правилам транслитерации записываются двумя буквами латинского алфавита (ю-yu, ё-yo и
т.д.).
Такие инициалы \verb+biblatex+ будет сокращать до одной буквы, что неверно.
Поправить его работу можно использовав опцию \verb+giveninits=false+.
Пример использования этой опции можно видеть в ссылке~\cite{initials}.

%\section{Формулы}\label{sec:ch1/sec3}

Благодаря пакету \textit{icomma}, \LaTeX~одинаково хорошо воспринимает
в~качестве десятичного разделителя и запятую (\(3,1415\)), и точку (\(3.1415\)).

\subsection{Ненумерованные одиночные формулы}\label{subsec:ch1/sec3/sub1}

Вот так может выглядеть формула, которую необходимо вставить в~строку
по~тексту: \(x \approx \sin x\) при \(x \to 0\).

А вот так выглядит ненумерованная отдельностоящая формула c подстрочными
и надстрочными индексами:
\[
    (x_1+x_2)^2 = x_1^2 + 2 x_1 x_2 + x_2^2
\]

Формула с неопределенным интегралом:
\[
    \int f(\alpha+x)=\sum\beta
\]

При использовании дробей формулы могут получаться очень высокие:
\[
    \frac{1}{\sqrt{2}+
        \displaystyle\frac{1}{\sqrt{2}+
            \displaystyle\frac{1}{\sqrt{2}+\cdots}}}
\]

В формулах можно использовать греческие буквы:
%Все \original... команды заранее, ради этого примера, определены в Dissertation\userstyles.tex
\[
    \alpha\beta\gamma\delta\originalepsilon\epsilon\zeta\eta\theta%
    \vartheta\iota\kappa\varkappa\lambda\mu\nu\xi\pi\varpi\rho\varrho%
    \sigma\varsigma\tau\upsilon\originalphi\phi\chi\psi\omega\Gamma\Delta%
    \Theta\Lambda\Xi\Pi\Sigma\Upsilon\Phi\Psi\Omega
\]
\[%https://texfaq.org/FAQ-boldgreek
    \boldsymbol{\alpha\beta\gamma\delta\originalepsilon\epsilon\zeta\eta%
        \theta\vartheta\iota\kappa\varkappa\lambda\mu\nu\xi\pi\varpi\rho%
        \varrho\sigma\varsigma\tau\upsilon\originalphi\phi\chi\psi\omega\Gamma%
        \Delta\Theta\Lambda\Xi\Pi\Sigma\Upsilon\Phi\Psi\Omega}
\]

Для добавления формул можно использовать пары \verb+$+\dots\verb+$+ и \verb+$$+\dots\verb+$$+,
но~они считаются устаревшими.
Лучше использовать их функциональные аналоги \verb+\(+\dots\verb+\)+ и \verb+\[+\dots\verb+\]+.

\subsection{Ненумерованные многострочные формулы}\label{subsec:ch1/sec3/sub2}

Вот так можно написать две формулы, не нумеруя их, чтобы знаки <<равно>> были
строго друг под другом:
\begin{align}
    f_W & =  \min \left( 1, \max \left( 0, \frac{W_{soil} / W_{max}}{W_{crit}} \right)  \right), \nonumber \\
    f_T & =  \min \left( 1, \max \left( 0, \frac{T_s / T_{melt}}{T_{crit}} \right)  \right), \nonumber
\end{align}

Выровнять систему ещё и по переменной \( x \) можно, используя окружение
\verb|alignedat| из пакета \verb|amsmath|. Вот так:
\[
|x| = \left\{
\begin{alignedat}{2}
     &   & x, \quad & \text{eсли } x\geqslant 0 \\
     & - & x, \quad & \text{eсли } x<0
\end{alignedat}
\right.
\]
Здесь первый амперсанд (в исходном \LaTeX\ описании формулы) означает
выравнивание по~левому краю, второй "--- по~\( x \), а~третий "--- по~слову
<<если>>. Команда \verb|\quad| делает большой горизонтальный пробел.

Ещё вариант:
\[
    |x|=
    \begin{cases}
        \phantom{-}x, \text{если } x \geqslant 0 \\
        -x, \text{если } x<0
    \end{cases}
\]

Кроме того, для  нумерованных формул \verb|alignedat| делает вертикальное
выравнивание номера формулы по центру формулы. Например, выравнивание
компонент вектора:
\begin{equation}
    \label{eq:2p3}
    \begin{alignedat}{2}
        {\mathbf{N}}_{o1n}^{(j)} = \,{\sin} \phi\,n\!\left(n+1\right)
        {\sin}\theta\,
        \pi_n\!\left({\cos} \theta\right)
        \frac{
        z_n^{(j)}\!\left( \rho \right)
        }{\rho}\,
         & {\boldsymbol{\hat{\mathrm e}}}_{r}\,+      \\
        +\,
        {\sin} \phi\,
        \tau_n\!\left({\cos} \theta\right)
        \frac{
        \left[\rho z_n^{(j)}\!\left( \rho \right)\right]^{\prime}
        }{\rho}\,
         & {\boldsymbol{\hat{\mathrm e}}}_{\theta}\,+ \\
        +\,
        {\cos} \phi\,
        \pi_n\!\left({\cos} \theta\right)
        \frac{
        \left[\rho z_n^{(j)}\!\left( \rho \right)\right]^{\prime}
        }{\rho}\,
         & {\boldsymbol{\hat{\mathrm e}}}_{\phi}\:.
    \end{alignedat}
\end{equation}

Ещё об отступах. Иногда для лучшей <<читаемости>> формул полезно
немного исправить стандартные интервалы \LaTeX\ с учётом логической
структуры самой формулы. Например в формуле~\cref{eq:2p3} добавлен
небольшой отступ \verb+\,+ между основными сомножителями, ниже
результат применения всех вариантов отступа:
\begin{align*}
    \backslash!             & \quad f(x) = x^2\! +3x\! +2         \\
    \mbox{по-умолчанию}     & \quad f(x) = x^2+3x+2               \\
    \backslash,             & \quad f(x) = x^2\, +3x\, +2         \\
    \backslash{:}           & \quad f(x) = x^2\: +3x\: +2         \\
    \backslash;             & \quad f(x) = x^2\; +3x\; +2         \\
    \backslash \mbox{space} & \quad f(x) = x^2\ +3x\ +2           \\
    \backslash \mbox{quad}  & \quad f(x) = x^2\quad +3x\quad +2   \\
    \backslash \mbox{qquad} & \quad f(x) = x^2\qquad +3x\qquad +2
\end{align*}

Можно использовать разные математические алфавиты:
\begin{align}
    \mathcal{ABCDEFGHIJKLMNOPQRSTUVWXYZ} \nonumber  \\
    \mathfrak{ABCDEFGHIJKLMNOPQRSTUVWXYZ} \nonumber \\
    \mathbb{ABCDEFGHIJKLMNOPQRSTUVWXYZ} \nonumber
\end{align}

Посмотрим на систему уравнений на примере аттрактора Лоренца:

\[
\left\{
\begin{array}{rl}
    \dot x = & \sigma (y-x)  \\
    \dot y = & x (r - z) - y \\
    \dot z = & xy - bz
\end{array}
\right.
\]

А для вёрстки матриц удобно использовать многоточия:
\[
    \left(
        \begin{array}{ccc}
            a_{11} & \ldots & a_{1n} \\
            \vdots & \ddots & \vdots \\
            a_{n1} & \ldots & a_{nn} \\
        \end{array}
    \right)
\]

\subsection{Нумерованные формулы}\label{subsec:ch1/sec3/sub3}

А вот так пишется нумерованная формула:
\begin{equation}
    \label{eq:equation1}
    e = \lim_{n \to \infty} \left( 1+\frac{1}{n} \right) ^n
\end{equation}

Нумерованных формул может быть несколько:
\begin{equation}
    \label{eq:equation2}
    \lim_{n \to \infty} \sum_{k=1}^n \frac{1}{k^2} = \frac{\pi^2}{6}
\end{equation}

Впоследствии на формулы~\cref{eq:equation1, eq:equation2} можно ссылаться.

Сделать так, чтобы номер формулы стоял напротив средней строки, можно,
используя окружение \verb|multlined| (пакет \verb|mathtools|) вместо
\verb|multline| внутри окружения \verb|equation|. Вот так:
\begin{equation} % \tag{S} % tag - вписывает свой текст
    \label{eq:equation3}
    \begin{multlined}
        1+ 2+3+4+5+6+7+\dots + \\
        + 50+51+52+53+54+55+56+57 + \dots + \\
        + 96+97+98+99+100=5050
    \end{multlined}
\end{equation}

Уравнения~\cref{eq:subeq_1,eq:subeq_2} демонстрируют возможности
окружения \verb|subequations| (пакет \verb|amsmath|).
\begin{subequations}
    \label{eq:subeq_1}
    \begin{gather}
        y = x^2 + 1 \label{eq:subeq_1-1} \\
        y = 2 x^2 - x + 1 \label{eq:subeq_1-2}
    \end{gather}
\end{subequations}
Ссылки на отдельные уравнения~\cref{eq:subeq_1-1,eq:subeq_1-2,eq:subeq_2-1}.
\begin{subequations}
    \label{eq:subeq_2}
    \begin{align}
        y & = x^3 + x^2 + x + 1 \label{eq:subeq_2-1} \\
        y & = x^2
    \end{align}
\end{subequations}

\subsection{Форматирование чисел и размерностей величин}\label{sec:units}

Числа форматируются при помощи команды \verb|\num|:
\num{5,3};
\num{2,3e8};
\num{12345,67890};
\num{2,6 d4};
\num{1+-2i};
\num{.3e45};
\num[exponent-base=2]{5 e64};
\num[exponent-base=2,exponent-to-prefix]{5 e64};
\num{1.654 x 2.34 x 3.430}
\num{1 2 x 3 / 4}.
Для написания последовательности чисел можно использовать команды \verb|\numlist| и \verb|\numrange|:
\numlist{10;30;50;70}; \numrange{10}{30}.
Значения углов можно форматировать при помощи команды \verb|\ang|:
\ang{2.67};
\ang{30,3};
\ang{-1;;};
\ang{;-2;};
\ang{;;-3};
\ang{300;10;1}.

Обратите внимание, что ГОСТ запрещает использование знака <<->> для обозначения отрицательных чисел
за исключением формул, таблиц и~рисунков.
Вместо него следует использовать слово <<минус>>.

Размерности можно записывать при помощи команд \verb|\si| и \verb|\SI|:
\si{\farad\squared\lumen\candela};
\si{\joule\per\mole\per\kelvin};
\si[per-mode = symbol-or-fraction]{\joule\per\mole\per\kelvin};
\si{\metre\per\second\squared};
\SI{0.10(5)}{\neper};
\SI{1.2-3i e5}{\joule\per\mole\per\kelvin};
\SIlist{1;2;3;4}{\tesla};
\SIrange{50}{100}{\volt}.
Список единиц измерений приведён в таблицах~\cref{tab:unit:base,
    tab:unit:derived,tab:unit:accepted,tab:unit:physical,tab:unit:other}.
Приставки единиц приведены в~таблице~\cref{tab:unit:prefix}.

С дополнительными опциями форматирования можно ознакомиться в~описании пакета \texttt{siunitx};
изменить или добавить единицы измерений можно в~файле \texttt{siunitx.cfg}.

\begin{table}
    \centering
    \captionsetup{justification=centering} % выравнивание подписи по-центру
    \caption{Основные величины СИ}\label{tab:unit:base}
    \begin{tabular}{llc}
        \toprule
        Название  & Команда          & Символ         \\
        \midrule
        Ампер     & \verb|\ampere|   & \si{\ampere}   \\
        Кандела   & \verb|\candela|  & \si{\candela}  \\
        Кельвин   & \verb|\kelvin|   & \si{\kelvin}   \\
        Килограмм & \verb|\kilogram| & \si{\kilogram} \\
        Метр      & \verb|\metre|    & \si{\metre}    \\
        Моль      & \verb|\mole|     & \si{\mole}     \\
        Секунда   & \verb|\second|   & \si{\second}   \\
        \bottomrule
    \end{tabular}
\end{table}

\begin{table}
    \small
    \centering
    \begin{threeparttable}% выравнивание подписи по границам таблицы
        \caption{Производные единицы СИ}\label{tab:unit:derived}
        \begin{tabular}{llc|llc}
            \toprule
            Название       & Команда               & Символ              & Название & Команда & Символ \\
            \midrule
            Беккерель      & \verb|\becquerel|     & \si{\becquerel}     &
            Ньютон         & \verb|\newton|        & \si{\newton}                                      \\
            Градус Цельсия & \verb|\degreeCelsius| & \si{\degreeCelsius} &
            Ом             & \verb|\ohm|           & \si{\ohm}                                         \\
            Кулон          & \verb|\coulomb|       & \si{\coulomb}       &
            Паскаль        & \verb|\pascal|        & \si{\pascal}                                      \\
            Фарад          & \verb|\farad|         & \si{\farad}         &
            Радиан         & \verb|\radian|        & \si{\radian}                                      \\
            Грей           & \verb|\gray|          & \si{\gray}          &
            Сименс         & \verb|\siemens|       & \si{\siemens}                                     \\
            Герц           & \verb|\hertz|         & \si{\hertz}         &
            Зиверт         & \verb|\sievert|       & \si{\sievert}                                     \\
            Генри          & \verb|\henry|         & \si{\henry}         &
            Стерадиан      & \verb|\steradian|     & \si{\steradian}                                   \\
            Джоуль         & \verb|\joule|         & \si{\joule}         &
            Тесла          & \verb|\tesla|         & \si{\tesla}                                       \\
            Катал          & \verb|\katal|         & \si{\katal}         &
            Вольт          & \verb|\volt|          & \si{\volt}                                        \\
            Люмен          & \verb|\lumen|         & \si{\lumen}         &
            Ватт           & \verb|\watt|          & \si{\watt}                                        \\
            Люкс           & \verb|\lux|           & \si{\lux}           &
            Вебер          & \verb|\weber|         & \si{\weber}                                       \\
            \bottomrule
        \end{tabular}
    \end{threeparttable}
\end{table}

\begin{table}
    \centering
    \begin{threeparttable}% выравнивание подписи по границам таблицы
        \caption{Внесистемные единицы}\label{tab:unit:accepted}

        \begin{tabular}{llc}
            \toprule
            Название        & Команда           & Символ          \\
            \midrule
            День            & \verb|\day|       & \si{\day}       \\
            Градус          & \verb|\degree|    & \si{\degree}    \\
            Гектар          & \verb|\hectare|   & \si{\hectare}   \\
            Час             & \verb|\hour|      & \si{\hour}      \\
            Литр            & \verb|\litre|     & \si{\litre}     \\
            Угловая минута  & \verb|\arcminute| & \si{\arcminute} \\
            Угловая секунда & \verb|\arcsecond| & \si{\arcsecond} \\ %
            Минута          & \verb|\minute|    & \si{\minute}    \\
            Тонна           & \verb|\tonne|     & \si{\tonne}     \\
            \bottomrule
        \end{tabular}
    \end{threeparttable}
\end{table}

\begin{table}
    \centering
    \captionsetup{justification=centering}
    \caption{Внесистемные единицы, получаемые из эксперимента}\label{tab:unit:physical}
    \begin{tabular}{llc}
        \toprule
        Название                & Команда                  & Символ                 \\
        \midrule
        Астрономическая единица & \verb|\astronomicalunit| & \si{\astronomicalunit} \\
        Атомная единица массы   & \verb|\atomicmassunit|   & \si{\atomicmassunit}   \\
        Боровский радиус        & \verb|\bohr|             & \si{\bohr}             \\
        Скорость света          & \verb|\clight|           & \si{\clight}           \\
        Дальтон                 & \verb|\dalton|           & \si{\dalton}           \\
        Масса электрона         & \verb|\electronmass|     & \si{\electronmass}     \\
        Электрон Вольт          & \verb|\electronvolt|     & \si{\electronvolt}     \\
        Элементарный заряд      & \verb|\elementarycharge| & \si{\elementarycharge} \\
        Энергия Хартри          & \verb|\hartree|          & \si{\hartree}          \\
        Постоянная Планка       & \verb|\planckbar|        & \si{\planckbar}        \\
        \bottomrule
    \end{tabular}
\end{table}

\begin{table}
    \centering
    \begin{threeparttable}% выравнивание подписи по границам таблицы
        \caption{Другие внесистемные единицы}\label{tab:unit:other}
        \begin{tabular}{llc}
            \toprule
            Название                  & Команда              & Символ             \\
            \midrule
            Ангстрем                  & \verb|\angstrom|     & \si{\angstrom}     \\
            Бар                       & \verb|\bar|          & \si{\bar}          \\
            Барн                      & \verb|\barn|         & \si{\barn}         \\
            Бел                       & \verb|\bel|          & \si{\bel}          \\
            Децибел                   & \verb|\decibel|      & \si{\decibel}      \\
            Узел                      & \verb|\knot|         & \si{\knot}         \\
            Миллиметр ртутного столба & \verb|\mmHg|         & \si{\mmHg}         \\
            Морская миля              & \verb|\nauticalmile| & \si{\nauticalmile} \\
            Непер                     & \verb|\neper|        & \si{\neper}        \\
            \bottomrule
        \end{tabular}
    \end{threeparttable}
\end{table}

\begin{table}
    \small
    \centering
    \begin{threeparttable}% выравнивание подписи по границам таблицы
        \caption{Приставки СИ}\label{tab:unit:prefix}
        \begin{tabular}{llcc|llcc}
            \toprule
            Приставка & Команда       & Символ      & Степень      &
            Приставка & Команда       & Символ      & Степень        \\
            \midrule
            Иокто     & \verb|\yocto| & \si{\yocto} & \textminus24 &
            Дека      & \verb|\deca|  & \si{\deca}  & 1              \\
            Зепто     & \verb|\zepto| & \si{\zepto} & \textminus21 &
            Гекто     & \verb|\hecto| & \si{\hecto} & 2              \\
            Атто      & \verb|\atto|  & \si{\atto}  & \textminus18 &
            Кило      & \verb|\kilo|  & \si{\kilo}  & 3              \\
            Фемто     & \verb|\femto| & \si{\femto} & \textminus15 &
            Мега      & \verb|\mega|  & \si{\mega}  & 6              \\
            Пико      & \verb|\pico|  & \si{\pico}  & \textminus12 &
            Гига      & \verb|\giga|  & \si{\giga}  & 9              \\
            Нано      & \verb|\nano|  & \si{\nano}  & \textminus9  &
            Терра     & \verb|\tera|  & \si{\tera}  & 12             \\
            Микро     & \verb|\micro| & \si{\micro} & \textminus6  &
            Пета      & \verb|\peta|  & \si{\peta}  & 15             \\
            Милли     & \verb|\milli| & \si{\milli} & \textminus3  &
            Екса      & \verb|\exa|   & \si{\exa}   & 18             \\
            Санти     & \verb|\centi| & \si{\centi} & \textminus2  &
            Зетта     & \verb|\zetta| & \si{\zetta} & 21             \\
            Деци      & \verb|\deci|  & \si{\deci}  & \textminus1  &
            Иотта     & \verb|\yotta| & \si{\yotta} & 24             \\
            \bottomrule
        \end{tabular}
    \end{threeparttable}
\end{table}

\subsection{Заголовки с формулами: \texorpdfstring{\(a^2 + b^2 = c^2\)}{%
        a\texttwosuperior\ + b\texttwosuperior\ = c\texttwosuperior},
    \texorpdfstring{\(\left\vert\textrm{{Im}}\Sigma\left(
            \protect\varepsilon\right)\right\vert\approx const\)}{|ImΣ (ε)| ≈ const},
    \texorpdfstring{\(\sigma_{xx}^{(1)}\)}{σ\_\{xx\}\textasciicircum\{(1)\}}
}\label{subsec:with_math}

Пакет \texttt{hyperref} берёт текст для закладок в pdf-файле из~аргументов
команд типа \verb|\section|, которые могут содержать математические формулы,
а~также изменения цвета текста или шрифта, которые не отображаются в~закладках.
Чтобы использование формул в заголовках не вызывало в~логе компиляции появление
предупреждений типа <<\texttt{Token not allowed in~a~PDF string
    (Unicode):(hyperref) removing...}>>, следует использовать конструкцию
\verb|\texorpdfstring{}{}|, где в~первых фигурных скобках указывается
формула, а~во~вторых "--- запись формулы для закладок.

\section{Рецензирование текста}\label{sec:markup}

В шаблоне для диссертации и автореферата заданы команды рецензирования.
Они видны при компиляции шаблона в режиме черновика или при установке
соответствующей настройки (\verb+showmarkup+) в~файле \verb+common/setup.tex+.

Команда \verb+\todo+ отмечает текст красным цветом.
\todo{Например, так.}

Команда \verb+\note+ позволяет выбрать цвет текста.
\note{Чёрный, } \note[red]{красный, } \note[green]{зелёный, }
\note[blue]{синий.} \note[orange]{Обратите внимание на ширину и расстановку
    формирующихся пробелов, в~результате приведённой записи (зависит также
    от~применяемого компилятора).}

Окружение \verb+commentbox+ также позволяет выбрать цвет.

\begin{commentbox}[red]
    Красный текст.

    Несколько параграфов красного текста.
\end{commentbox}

\begin{commentbox}[blue]
    Синяя формула.

    \begin{equation}
        \alpha + \beta = \gamma
    \end{equation}
\end{commentbox}

\verb+commentbox+ позволяет закомментировать участок кода в~режиме чистовика.
Чтобы убрать кусок кода для всех режимов, можно использовать окружение
\verb+comment+.

\begin{comment}
Этот текст всегда скрыт.
\end{comment}

\section{Работа со списком сокращений и~условных обозначений}\label{sec:acronyms}

С помощью пакета \texttt{nomencl} можно создавать удобный сортированный список
сокращений и условных обозначений во время написания текста. Вызов
\verb+\nomenclature+ добавляет нужный символ или сокращение с~описанием
в~список, который затем печатается вызовом \verb+\printnomenclature+
в~соответствующем разделе.
Для того, чтобы эти операции прошли, потребуется дополнительный вызов
\verb+makeindex -s nomencl.ist -o %.nls %.nlo+ в~командной строке, где вместо
\verb+%+ следует подставить имя главного файла проекта (\verb+dissertation+
для этого шаблона).
Затем потребуется один или два дополнительных вызова компилятора проекта.
\begin{equation}
    \omega = c k,
\end{equation}
где \( \omega \) "--- частота света, \( c \) "--- скорость света, \( k \) "---
модуль волнового вектора.
\nomenclature{\(\omega\)}{частота света\nomrefeq}
\nomenclature{\(c\)}{скорость света\nomrefpage}
\nomenclature{\(k\)}{модуль волнового вектора\nomrefeqpage}
Использование
\begin{verbatim}
\nomenclature{\(\omega\)}{частота света\nomrefeq}
\nomenclature{\(c\)}{скорость света\nomrefpage}
\nomenclature{\(k\)}{модуль волнового вектора\nomrefeqpage}
\end{verbatim}
после уравнения добавит в список условных обозначений три записи.
Ссылки \verb+\nomrefeq+ на последнее уравнение, \verb+\nomrefpage+ "--- на
страницу, \verb+\nomrefeqpage+ "--- сразу на~последнее уравнение и~на~страницу,
можно опускать и~не~использовать.

Группировкой и сортировкой пунктов в списке можно управлять с~помощью указания
дополнительных аргументов к команде \verb+nomenclature+.
Например, при вызове
\begin{verbatim}
\nomenclature[03]{\( \hbar \)}{постоянная Планка}
\nomenclature[01]{\( G \)}{гравитационная постоянная}
\end{verbatim}
\( G \) будет стоять в списке выше, чем \( \hbar \).
Для корректных вертикальных отступов между строками в описании лучше
не~использовать многострочные формулы в~списке обозначений.

\nomenclature{%
    \( \begin{rcases}
        a_n \\
        b_n
    \end{rcases} \)%
}{коэффициенты разложения Ми в дальнем поле соответствующие электрическим и
    магнитным мультиполям}
\nomenclature[a\( e \)]{\( {\boldsymbol{\hat{\mathrm e}}} \)}{единичный вектор}
\nomenclature{\( E_0 \)}{амплитуда падающего поля}
\nomenclature{\( j \)}{тип функции Бесселя}
\nomenclature{\( k \)}{волновой вектор падающей волны}
\nomenclature{%
    \( \begin{rcases}
        a_n \\
        b_n
    \end{rcases} \)%
}{и снова коэффициенты разложения Ми в дальнем поле соответствующие
    электрическим и магнитным мультиполям. Добавлено много текста, так что
    описание группы условных обозначений значительно превысило высоту этой
    группы...}
\nomenclature{\( L \)}{общее число слоёв}
\nomenclature{\( l \)}{номер слоя внутри стратифицированной сферы}
\nomenclature{\( \lambda \)}{длина волны электромагнитного излучения в вакууме}
\nomenclature{\( n \)}{порядок мультиполя}
\nomenclature{%
    \( \begin{rcases}
        {\mathbf{N}}_{e1n}^{(j)} & {\mathbf{N}}_{o1n}^{(j)} \\
        {\mathbf{M}_{o1n}^{(j)}} & {\mathbf{M}_{e1n}^{(j)}}
    \end{rcases} \)%
}{сферические векторные гармоники}
\nomenclature{\( \mu \)}{магнитная проницаемость в вакууме}
\nomenclature{\( r, \theta, \phi \)}{полярные координаты}
\nomenclature{\( \omega \)}{частота падающей волны}

С помощью \verb+nomenclature+ можно включать в~список сокращения,
не~используя их~в~тексте.
% запись сокращения в список происходит командой \nomenclature,
% а не употреблением самого сокращения
\nomenclature{FEM}{finite element method, метод конечных элементов}
\nomenclature{FIT}{finite integration technique, метод конечных интегралов}
\nomenclature{FMM}{fast multipole method, быстрый метод многополюсника}
\nomenclature{FVTD}{finite volume time-domain, метод конечных объёмов
    во~временной области}
\nomenclature{MLFMA}{multilevel fast multipole algorithm, многоуровневый
    быстрый алгоритм многополюсника}
\nomenclature{BEM}{boundary element method, метод граничных элементов}
\nomenclature{CST MWS}{Computer Simulation Technology Microwave Studio
    программа для компьютерного моделирования уравнен Максвелла}
\nomenclature{DDA}{discrete dipole approximation, приближение дискретиных
    диполей}
\nomenclature{FDFD}{finite difference frequency domain, метод конечных
    разностей в~частотной области}
\nomenclature{FDTD}{finite difference time domain, метод конечных разностей
    во~временной области}
\nomenclature{MoM}{method of moments, метод моментов}
\nomenclature{MSTM}{multiple sphere T-Matrix, метод Т-матриц для множества
    сфер}
\nomenclature{PSTD}{pseudospectral time domain method, псевдоспектральный метод
    во~временной области}
\nomenclature{TLM}{transmission line matrix method, метод матриц линий передач}

\FloatBarrier
